\part{Proofs of the arithmetic inputs}\label{part:proofs}

\section{Introduction to Part II}\label{sec:pf:intro}

Part~I proves Theorems~A and~B and Corollary~C from the arithmetic inputs stated in
\cref{sec:top,sec:stepanov,sec:families,sec:window,sec:galois,sec:covolume}, and outlines their proofs.
Part~II gives the proofs. They are the proofs of the formal development \cite{LeanSource2026}, whose
verification is described in \cref{sec:verification}, written out in mathematical language section by section.
Each section of Part~II treats one part of the formal proof: it restates the results of Part~I that it
proves, introduces the notions it needs, and proves its results from lemmas that are stated precisely.

\subsection{Conventions}\label{sec:pf:intro:conv}
\begin{enumerate}[label=(\arabic*)]
\item A result that corresponds to a declaration of the formal development carries the name of the
declaration in brackets at the end of its statement, as in \formal{UpperTheory.Zero.LLC.recTExists}.
Several names are attached to one statement when the development contains several equivalent forms of it.
\item A reference to a result of the formal development is printed as a cross-reference (``Lemma~23.4'')
when the result is proved in this paper, and as the bracketed name of the declaration otherwise. Results of
the second kind are called \emph{deferred}: the statement used is paraphrased where it is cited, and its
proof is that of the formal development. They are the deepest lemmas of each part of the proof (see
\cref{sec:pf:intro:scope}).
\item Results from the literature are quoted with references. When a formal statement differs from the
published one, a remark says how.
\item The formal development uses conventions that are invisible in ordinary mathematical writing: division
by $0$ gives $0$; the length of a module of infinite length and the cardinality of an infinite set are $0$;
a divergent sum is $0$; the supremum of an unbounded set of reals is $0$. When a statement depends on such a
convention, this is pointed out.
\item Remarks record hypotheses of formal statements that are not used in their proofs, statements that hold
for trivial reasons, and names of formal declarations that do not describe their content. These remarks
have no mathematical consequences; they are included to make the correspondence with the formal
development exact.
\end{enumerate}

\subsection{Scope}\label{sec:pf:intro:scope}
The formal proofs of the arithmetic inputs consist of about $4.8$ million lines. Part~II converts each part
of them down to a certain depth. Every result of Part~I that is outlined there is proved here from lemmas
that are stated precisely, and most of these lemmas are proved in turn, usually through two or three
further levels. In all, Part~II states and proves $6\,770$ declarations of the formal development. The
deeper lemmas whose proofs are not reproduced ($7\,295$ declarations, some of them with long formal proofs)
are cited by the names of their formal counterparts, as explained in (2) above. Classical
theorems that the development formalizes are cited, not reproved.

\subsection{Organisation}\label{sec:pf:intro:org}
\Cref{tab:pf:map} lists, for each result of Part~I that is proved in Part~II, the section containing its proof.
The sections are grouped as follows.
\begin{itemize}
\item The top-level deduction and the Stepanov-type bound (Sections 6 and 7): \cref{sec:pf:top-supply}, \cref{sec:pf:top-shape}, \cref{sec:pf:linkG}.
\item The hypothesis package (Section 8): \cref{sec:pf:stageU}, \cref{sec:pf:liftinj}, \cref{sec:pf:bundles}, \cref{sec:pf:baselat2}, \cref{sec:pf:baselat1}, \cref{sec:pf:rows}, \cref{sec:pf:reckernel}, \cref{sec:pf:thmS}, \cref{sec:pf:thmS-leaves}, \cref{sec:pf:thmS-carrier}, \cref{sec:pf:thmK}, \cref{sec:pf:linkN}, \cref{sec:pf:linkN-sel}, \cref{sec:pf:linkN-adm}, \cref{sec:pf:onto3}, \cref{sec:pf:onto1}, \cref{sec:pf:onto2}, \cref{sec:pf:open0}, \cref{sec:pf:open0a}, \cref{sec:pf:rowprops}, \cref{sec:pf:fit}, \cref{sec:pf:fs}, \cref{sec:pf:multnc}.
\item The window bound (Section 9): \cref{sec:pf:vwin}, \cref{sec:pf:vwin-he}, \cref{sec:pf:kato}, \cref{sec:pf:eis}, \cref{sec:pf:modularity}.
\item Galois representations and local parameters (Section 10): \cref{sec:pf:llc}, \cref{sec:pf:geo}, \cref{sec:pf:hlf-cech}, \cref{sec:pf:hlf-chart}, \cref{sec:pf:galred}, \cref{sec:pf:tlift}, \cref{sec:pf:steps2}, \cref{sec:pf:steps1}, \cref{sec:pf:ks2}, \cref{sec:pf:ks1}.
\item Covolume bounds (Section 11): \cref{sec:pf:covol}, \cref{sec:pf:lift2}, \cref{sec:pf:lift1}, \cref{sec:pf:deepstr}, \cref{sec:pf:tower}, \cref{sec:pf:towergeo}, \cref{sec:pf:volthm}.
\item Results used in several parts of the proof: \cref{sec:pf:sh-theta}, \cref{sec:pf:sh-flag}, \cref{sec:pf:sh-siegel}, \cref{sec:pf:sh-div}, \cref{sec:pf:sh-hodge}, \cref{sec:pf:sh-curv}, \cref{sec:pf:sh-local}, \cref{sec:pf:sh-misc}.
\end{itemize}

{\small
\begin{longtable}{>{\raggedright\arraybackslash}p{0.3\textwidth}>{\raggedright\arraybackslash}p{0.6\textwidth}}
\caption{Results of Part~I and the sections of Part~II that prove them.}\label{tab:pf:map}\\
\toprule
result of Part I & proved in\\
\midrule
\endfirsthead
\toprule
result of Part I & proved in\\
\midrule
\endhead
\Cref{thm:supply} & \cref{sec:pf:top-supply} \\
\Cref{thm:top-deduction} & \cref{sec:pf:top-supply} \\
\Cref{thm:shape} & \cref{sec:pf:top-shape} \\
\Cref{thm:global-chain} & \cref{sec:pf:top-shape} \\
\Cref{thm:linkIRR} & \cref{sec:pf:top-shape} \\
\Cref{prop:families:admLq} & \cref{sec:pf:top-shape} \\
\Cref{prop:families:div} & \cref{sec:pf:top-shape} \\
\Cref{cor:families:linkS} & \cref{sec:pf:top-shape} \\
\Cref{prop:families:famgl} & \cref{sec:pf:top-shape} \\
\Cref{prop:families:design} & \cref{sec:pf:top-shape} \\
\Cref{thm:linkG} & \cref{sec:pf:linkG} \\
\Cref{thm:families:local} & \cref{sec:pf:linkG} \\
\Cref{thm:hyps-nonempty} & \cref{sec:pf:stageU} \\
\Cref{prop:families:liftinj} & \cref{sec:pf:liftinj} \\
\Cref{thm:member-bundles} & \cref{sec:pf:bundles} \\
\Cref{prop:families:mass} & \cref{sec:pf:bundles} \\
\Cref{prop:families:allrows} & \cref{sec:pf:rows} \\
\Cref{thm:counting} & \cref{sec:pf:rows} \\
\Cref{prop:families:volratio} & \cref{sec:pf:rows} \\
\Cref{prop:families:periods} & \cref{sec:pf:rows} \\
\Cref{thm:theorem-S} & \cref{sec:pf:thmS} \\
\Cref{thm:theorem-K} & \cref{sec:pf:thmK} \\
\Cref{thm:linkN} & \cref{sec:pf:linkN} \\
\Cref{prop:families:points} & \cref{sec:pf:onto3}, \cref{sec:pf:fit} \\
\Cref{thm:families:open} & \cref{sec:pf:open0} \\
\Cref{prop:families:rowrat} & \cref{sec:pf:rowprops} \\
\Cref{prop:families:fs} & \cref{sec:pf:fs} \\
\Cref{thm:vwin} & \cref{sec:pf:vwin} \\
\Cref{prop:window:assembly} & \cref{sec:pf:vwin} \\
\Cref{prop:psibound} & \cref{sec:pf:vwin} \\
\Cref{lem:window:rec} & \cref{sec:pf:vwin} \\
\Cref{thm:window:lemmaK} & \cref{sec:pf:vwin} \\
\Cref{lem:window:tower} & \cref{sec:pf:vwin} \\
\Cref{thm:window:bab} & \cref{sec:pf:vwin} \\
\Cref{prop:window:propZ} & \cref{sec:pf:vwin} \\
\Cref{thm:plus-part} & \cref{sec:pf:vwin} \\
\Cref{thm:window:HE} & \cref{sec:pf:vwin-he} \\
\Cref{thm:kato} & \cref{sec:pf:kato} \\
\Cref{lem:window:heights} & \cref{sec:pf:kato} \\
\Cref{thm:window:countI} & \cref{sec:pf:eis} \\
\Cref{thm:rectexists} & \cref{sec:pf:llc} \\
\Cref{prop:galois:unique} & \cref{sec:pf:llc} \\
\Cref{thm:row0b} & \cref{sec:pf:geo} \\
\Cref{thm:geoone} & \cref{sec:pf:geo} \\
\Cref{thm:galois:hlf} & \cref{sec:pf:geo} \\
\Cref{thm:galois:reduction} & \cref{sec:pf:galred} \\
\Cref{thm:galois:tlift} & \cref{sec:pf:tlift} \\
\Cref{thm:covol} & \cref{sec:pf:covol} \\
\Cref{thm:covolume:assembly} & \cref{sec:pf:covol}, \cref{sec:pf:lift2} \\
\Cref{thm:covolume:lift} & \cref{sec:pf:lift2} \\
\Cref{thm:deepstr} & \cref{sec:pf:deepstr} \\
\Cref{thm:tower} & \cref{sec:pf:tower} \\
\bottomrule
\end{longtable}
}
